\section{Заключение}

В ходе выполнения лабораторной работы были детально изучены и практически освоены принципы распределённой системы контроля версий \texttt{Git} и централизованной системы \texttt{Subversion}. Таким образом реализована имитация многопользовательской разработки, что позволило отработать ключевые аспекты командной работы над проектом.

Были закреплены следующие практические навыки:

\begin{itemize}
  \item создание и переключение между ветками разработки (\texttt{branch}, \texttt{switch}, \texttt{checkout});
  \item внесение изменений в рабочую копию с последующей фиксацией (\texttt{add}, \texttt{commit});
  \item синхронизация с удалёнными репозиториями (\texttt{push}, \texttt{fetch}, \texttt{update});
  \item копирование веток в централизованной модели (\texttt{copy}).
\end{itemize}

Особое внимание уделено разрешению конфликтов, которые возникают при параллельном изменении одних и тех же файлов разными участниками, что часто встречается в реальной командной разработке.